\documentclass[12pt,a4paper]{article}

\usepackage[margin=1in]{geometry}
\usepackage{booktabs}
\usepackage{array}
\usepackage{longtable}
\usepackage{hyperref}
\usepackage{xcolor}
\usepackage{enumitem}
\usepackage{float}
\usepackage{caption}

\hypersetup{
    colorlinks=true,
    linkcolor=blue,
    urlcolor=blue,
    citecolor=blue
}

\setlist[itemize]{leftmargin=1.2em}
\setlist[enumerate]{leftmargin=1.4em}
\renewcommand{\arraystretch}{1.18}

\title{\textbf{RACE Tutor AI}\\
\large Classical Machine Learning System for Reading Comprehension Support}
\author{
Abdullah Sajid\\
Roll No. 22i-2346
\and
Huzaifa Khalid\\
Roll No. 22i-2332
}
\date{May 2026}

\begin{document}

\maketitle

\begin{abstract}
RACE Tutor AI is an interactive reading-comprehension tutoring application built
on the RACE dataset. The system combines a classical machine-learning answer
verification model with an extractive distractor and hint generator, then exposes
the complete workflow through a Streamlit interface. Model A predicts the most
likely answer option for a given passage and question. Model B generates
plausible distractors and graduated hints from the source passage. The final
Model A configuration uses one candidate example per answer option, binary
bag-of-words features, bigrams, TF-IDF similarity, lexical overlap, negation
features, and per-question relative option features. On the held-out test set,
the final Model A reaches 37.81\% four-option answer-selection accuracy.
\end{abstract}

\tableofcontents
\newpage

\section{Introduction}

Reading-comprehension systems are useful for education because they can help
students practise passage analysis, answer selection, and evidence-based
reasoning. This project implements a complete tutoring pipeline for the RACE
dataset using transparent classical machine-learning methods. The goal is not to
build a large neural model, but to create a reproducible system that can load
RACE-style data, train an answer verifier, generate hints and distractors, and
run inside an interactive application.

The project has three main components:

\begin{itemize}
    \item \textbf{Model A}: an answer-verification and answer-selection model.
    \item \textbf{Model B}: a deterministic distractor and hint generator.
    \item \textbf{Frontend}: a Streamlit user interface for loading quizzes,
    answering questions, revealing hints, and viewing session analytics.
\end{itemize}

\section{Dataset}

The project uses the RACE dataset, a large-scale reading-comprehension dataset
based on English examinations. Each example contains an article, a question,
four answer options, and the correct option label. The local project expects the
following flat CSV columns:

\begin{center}
\texttt{id, article, question, A, B, C, D, answer}
\end{center}

The preprocessing code can also load the HuggingFace \texttt{ehovy/race}
dataset and flatten each example from a nested \texttt{options} list into the
explicit \texttt{A}, \texttt{B}, \texttt{C}, and \texttt{D} columns. Examples
without exactly four options are ignored so that all downstream models receive a
consistent four-choice structure.

\subsection{Dataset Splits}

Table~\ref{tab:splits} shows the cleaned split sizes used during the final
training workflow. Cleaning removes invalid answers and rows with empty article
or question text.

\begin{table}[H]
\centering
\caption{Cleaned RACE split sizes used by the project.}
\label{tab:splits}
\begin{tabular}{lrrrr}
\toprule
\textbf{Split} & \textbf{Rows} & \textbf{Candidate Rows} & \textbf{Answer A} & \textbf{Not A} \\
\midrule
Train & 70,274 & 281,096 & 15,315 & 54,959 \\
Validation & 8,784 & 35,136 & 1,914 & 6,870 \\
Test & 8,787 & 35,148 & 1,915 & 6,872 \\
\bottomrule
\end{tabular}
\end{table}

\section{Preprocessing}

The preprocessing pipeline is implemented in \path{backend/preprocessing.py}
and \path{backend/feature_engineering.py}. It performs the following steps:

\begin{enumerate}
    \item Load local RACE CSV files or flatten HuggingFace RACE examples.
    \item Keep only examples with four options.
    \item Normalize text by lowercasing, removing punctuation, filling missing
    text, and normalizing whitespace.
    \item Drop rows with invalid answer labels.
    \item Drop rows with empty articles or questions.
    \item Convert each question into candidate-level training examples.
\end{enumerate}

For the final Model A, each original question is expanded into four candidate
examples. Each candidate asks whether a particular option is the correct answer.
This framing matches the final answer-selection task better than training on
only option \texttt{A}.

\section{Model A: Answer Verification}

\subsection{Initial Framing}

The first Model A version used a binary option-\texttt{A} formulation:

\begin{center}
\textit{Is option A correct?}
\end{center}

The label was 1 when \texttt{answer == "A"} and 0 otherwise. This version was
useful as a baseline and produced preprocessed artifacts in
\texttt{data/processed/model\_a}. However, it used only one option per question,
which discarded information from options \texttt{B}, \texttt{C}, and \texttt{D}.

\subsection{Final Candidate-Level Framing}

The final model trains one example per candidate answer. For each RACE row, the
system creates four examples:

\begin{center}
\texttt{article + question + option A},\quad
\texttt{article + question + option B},\quad
\texttt{article + question + option C},\quad
\texttt{article + question + option D}
\end{center}

The label is 1 for the correct candidate and 0 for the other three candidates.
At inference time, the model scores all four options and selects the option with
the highest score.

\subsection{Feature Engineering}

The final feature set combines sparse text features and dense handcrafted
features:

\begin{itemize}
    \item Binary bag-of-words features from article, question, and candidate
    option text.
    \item Unigram and bigram features using \texttt{ngram\_max = 2}.
    \item TF-IDF cosine similarity between the article and
    \texttt{question + option}.
    \item TF-IDF cosine similarity between the article and question.
    \item TF-IDF cosine similarity between the article and candidate option.
    \item Article length, question length, and option length.
    \item Question-article, option-article, and option-question lexical overlap.
    \item Option coverage ratio in the article.
    \item Negation flags for question and option text.
    \item Per-question relative option features, including option length ratio
    and overlap deltas against the sibling options.
\end{itemize}

The vectorizers and dense-feature scaler are fitted only on the training split.
Validation and test data are transformed using the fitted training objects to
avoid data leakage.

\subsection{Training Configuration}

The best final Model A was trained with logistic regression:

\begin{table}[H]
\centering
\caption{Final Model A training configuration.}
\label{tab:model-a-config}
\begin{tabular}{ll}
\toprule
\textbf{Setting} & \textbf{Value} \\
\midrule
Model & Logistic Regression \\
Training formulation & Four candidate examples per question \\
Class weight & None \\
Regularization parameter & \texttt{C = 0.3} \\
Maximum vocabulary size & 50,000 \\
Minimum document frequency & 2 \\
Maximum n-gram size & 2 \\
Saved artifact & \texttt{models/model\_a/logreg\_candidate\_rich.pkl} \\
\bottomrule
\end{tabular}
\end{table}

The exact command used for final training was:

\begin{verbatim}
.venv/bin/python -m backend.train_model_a \
  --train data/raw/train.csv \
  --val data/raw/val.csv \
  --test data/raw/test.csv \
  --model logreg_candidate \
  --class-weight none \
  --ngram-max 2 \
  --c 0.3 \
  --output-name logreg_candidate_rich
\end{verbatim}

\subsection{Model A Results}

Table~\ref{tab:model-a-progression} shows the answer-selection improvement
during development.

\begin{table}[H]
\centering
\caption{Model A answer-selection accuracy progression.}
\label{tab:model-a-progression}
\begin{tabular}{lcc}
\toprule
\textbf{Model} & \textbf{Validation Accuracy} & \textbf{Test Accuracy} \\
\midrule
Option-A binary model & 0.3283 & 0.3214 \\
Candidate model & 0.3530 & 0.3570 \\
Rich candidate model & \textbf{0.3809} & \textbf{0.3781} \\
\bottomrule
\end{tabular}
\end{table}

The final binary candidate-level metrics are shown in
Table~\ref{tab:model-a-binary}. These metrics use the independent candidate
labels. The answer-selection metric is more important for the application
because the app ranks all four options for each question.

\begin{table}[H]
\centering
\caption{Final Model A binary candidate-level metrics.}
\label{tab:model-a-binary}
\begin{tabular}{lrrrr}
\toprule
\textbf{Split} & \textbf{Accuracy} & \textbf{Macro F1} & \textbf{Precision} & \textbf{Recall} \\
\midrule
Validation & 0.7446 & 0.4671 & 0.4054 & 0.0459 \\
Test & 0.7444 & 0.4660 & 0.4002 & 0.0447 \\
\bottomrule
\end{tabular}
\end{table}

Table~\ref{tab:model-a-answer} shows the final four-option answer-selection
metrics.

\begin{table}[H]
\centering
\caption{Final Model A four-option answer-selection metrics.}
\label{tab:model-a-answer}
\begin{tabular}{lcc}
\toprule
\textbf{Split} & \textbf{Answer Accuracy} & \textbf{Macro F1} \\
\midrule
Validation & 0.3809 & 0.3798 \\
Test & 0.3781 & 0.3771 \\
\bottomrule
\end{tabular}
\end{table}

\section{Model B: Distractor and Hint Generation}

Model B is implemented in \texttt{backend/model\_b.py}. It is a deterministic
classical component rather than a trained model artifact. This design keeps the
system lightweight and explainable.

\subsection{Distractor Generation}

The distractor generator extracts lightweight passage phrases without requiring
external NLP models. It:

\begin{enumerate}
    \item Splits the article into sentences.
    \item Extracts content words while removing English stopwords.
    \item Builds one- to four-word phrase candidates.
    \item Filters out phrases that duplicate or contain the correct answer.
    \item Scores candidates using passage frequency, question overlap, phrase
    length, and similarity to the correct answer.
    \item Returns three distinct distractors.
\end{enumerate}

\subsection{Hint Generation}

The hint generator ranks article sentences by overlap with the question and
correct answer. It returns three hints:

\begin{enumerate}
    \item A broad focus hint based on important question terms.
    \item A supporting sentence with answer terms masked.
    \item The strongest supporting sentence as direct evidence.
\end{enumerate}

This gives the student a gradual path from general guidance to stronger
evidence.

\section{Frontend Application}

The frontend is built in Streamlit and is implemented in
\texttt{frontend/app.py}. It provides four tabs:

\begin{itemize}
    \item \textbf{Article Input}: load a random RACE example or create a manual
    quiz.
    \item \textbf{Quiz View}: show the passage, question, answer options, and
    Model A prediction.
    \item \textbf{Hints}: reveal Model B hints one at a time and view generated
    distractors.
    \item \textbf{Analytics}: track attempts, accuracy, latency, and answer
    history for the current session.
\end{itemize}

The app uses \texttt{models/model\_a/logreg\_candidate\_rich.pkl} as its default
answer-verification model.

\section{Verification}

The completed project was verified through:

\begin{itemize}
    \item Python byte-compilation of backend, frontend, and test modules.
    \item Backend smoke test covering dataset loading, quiz construction, Model
    A loading, Model A prediction, Model B distractor generation, and Model B
    hint generation.
    \item Streamlit server startup.
    \item HTTP response check from the running Streamlit app.
\end{itemize}

The Streamlit app was started successfully at:

\begin{center}
\url{http://127.0.0.1:8501}
\end{center}

\section{Limitations}

Although the final model improves over the initial baselines, the system remains
limited by its classical feature representation. RACE questions often require
multi-sentence reasoning, inference, paraphrase recognition, and world
knowledge. Bag-of-words and overlap features cannot fully capture these
phenomena. Model B is also extractive and deterministic, so distractors can be
surface-level and may not always match the semantic type of the correct answer.

\section{Future Work}

Future improvements include:

\begin{itemize}
    \item Use stronger ranking models that optimize answer selection directly.
    \item Add semantic features from sentence embeddings or transformer models.
    \item Improve Model B with noun-phrase chunking and answer-type matching.
    \item Add persistent analytics and per-student progress tracking.
    \item Build a larger evaluation suite for generated distractor quality and
    hint usefulness.
\end{itemize}

\section{Conclusion}

RACE Tutor AI is now a complete running project. It loads and preprocesses RACE
data, trains a classical answer-selection model, generates distractors and
hints, and presents the workflow through a Streamlit interface. The strongest
final Model A reaches 37.81\% test answer-selection accuracy, improving over
the option-\texttt{A} baseline and the first candidate-level model. Model B
adds explainable educational support through passage-based distractors and
graduated hints.

\begin{thebibliography}{9}

\bibitem{race}
Guokun Lai, Qizhe Xie, Hanxiao Liu, Yiming Yang, and Eduard Hovy.
\textit{RACE: Large-scale ReAding Comprehension Dataset From Examinations}.
EMNLP, 2017.

\bibitem{sklearn}
Fabian Pedregosa et al.
\textit{Scikit-learn: Machine Learning in Python}.
Journal of Machine Learning Research, 2011.

\bibitem{streamlit}
Streamlit documentation.
\textit{Streamlit: A faster way to build and share data apps}.
\url{https://streamlit.io/}.

\end{thebibliography}

\end{document}
